% Section 4.1, from lectures/L04/appendix/a_factory_raw_pointers.md.

\section{A Simple Driver Factory with Raw Pointers}
\label{c4:sec:raw}\label{c4:app:a}
This section shows a simplified example of how a factory can be used to create drivers in an
embedded system.

The purpose is to demonstrate:
\begin{itemize}
  \item What a factory is.
  \item Why we use it.
  \item How it makes the system logic independent of hardware.
\end{itemize}

\subsection{Raw pointers and dynamic memory allocation}
\label{c4:sec:rawpointers}
In this example, so-called raw pointers are used to clarify what happens. Raw pointers are
obtained when we use the operators \code{new} and \code{delete} to allocate and free objects
respectively:
\begin{itemize}
  \item These operators serve a similar function to \code{malloc()} and \code{free()} in C, but in
    C++ \code{new} also calls the object's constructor and \code{delete} its destructor.
  \item Just as when using \code{malloc()} and \code{free()}, we are responsible for freeing
    allocated resources ourselves.
  \item In modern C++, smart pointers are used instead, which ensure that resources are released
    via the pointer's destructor when it goes out of scope, preventing memory leaks.
  \item Smart pointers are discussed in \secref{c4:sec:smart}, where these raw pointers are replaced
    with \code{std::unique_ptr}.
\end{itemize}

\subsection{Background}
\label{c4:sec:background}
Assume that we have implemented the following drivers as classes:
\begin{itemize}
  \item GPIO: \code{driver::gpio::Interface} with the following concrete implementations:
  \begin{itemize}
    \item \code{driver::gpio::Esp32s3}: GPIO driver for the ESP32-S3 microcontroller.
    \item \code{driver::gpio::Stub}: GPIO stub driver.
  \end{itemize}
\end{itemize}
Assume that we have created the following class \code{app::logic::Logic} to handle the system
logic:
\begin{itemize}
  \item The system logic uses an LED and a button, where the LED is toggled on the rising edge of
    the push button.
  \item Interfaces are used so that the hardware used can be selected generically, i.e.\ the
    concrete types are chosen by the user.
  \item The method \code{run()} is used to execute the system logic continuously.
\end{itemize}

\begin{cppcode}
namespace app::logic
{
class Logic final
{
public:
    explicit Logic(driver::gpio::Interface& led,
                   driver::gpio::Interface& button) noexcept
        : myLed{led}
        , myButton{button}
    {}

    ~Logic() noexcept
    {
        myLed.write(false);
    }

    void run() noexcept
    {
        bool buttonPrev{false};

        while (1)
        {
            const bool buttonCurrent{myButton.read()};
            if (buttonCurrent && !buttonPrev)
            {
                myLed.toggle();
            }
            buttonPrev = buttonCurrent;
        }
    }

    Logic()                        = delete;
    Logic(const Logic&)            = delete;
    Logic(Logic&&)                 = delete;
    Logic& operator=(const Logic&) = delete;
    Logic& operator=(Logic&&)      = delete;

private:
    driver::gpio::Interface& myLed;
    driver::gpio::Interface& myButton;
};
} // namespace app::logic
\end{cppcode}

\noindent The example above can be implemented by creating drivers of any concrete types before the
system logic is created, as shown below:

\begin{cppcode}
int main()
{
    constexpr std::uint8_t ledPin{2U};
    constexpr std::uint8_t buttonPin{3U};

    // Create ESP32-S3 drivers (ignoring data direction for simplicity).
    driver::gpio::Esp32s3 led{ledPin};
    driver::gpio::Esp32s3 button{buttonPin};

    // Create and run system logic.
    app::logic::Logic logic{led, button};
    logic.run();
}
\end{cppcode}

\noindent The method above is very efficient when dynamic memory allocation should be limited.

However, it is often desirable to create multiple drivers of the same type, for example drivers
for a specific microcontroller. For this reason, it is common to implement classes that create
such instances. Such a class is called a \textbf{factory}, and its sole purpose is to create
objects.

Assume that the following factories have been implemented:
\begin{itemize}
  \item \code{driver::factory::Interface} with the following concrete implementations:
  \begin{itemize}
    \item \code{driver::factory::Esp32s3}: Creates ESP32-S3 instances.
    \item \code{driver::factory::Stub}: Creates stub instances.
  \end{itemize}
\end{itemize}
Assume that the constructor for \code{app::logic::Logic} had taken a reference to an arbitrary
factory:

\begin{cppcode}
explicit Logic(driver::factory::Interface& factory,
               const std::uint8_t ledPin,
               const std::uint8_t buttonPin) noexcept;
\end{cppcode}

\noindent Assuming that the member variables had been implemented as pointers instead of
references:

\begin{cppcode}
driver::gpio::Interface* myLed;
driver::gpio::Interface* myButton;
\end{cppcode}

\noindent and the factory had created raw pointers, the instances could be created as shown below:

\begin{cppcode}
explicit Logic(driver::factory::Interface& factory,
               const std::uint8_t ledPin,
               const std::uint8_t buttonPin) noexcept
        : myLed{factory.gpio(ledPin)}
        , myButton{factory.gpio(buttonPin)}
    {}
\end{cppcode}

\noindent As an example, \code{Esp32s3} instances could be created by passing a
\code{driver::factory::Esp32s3} to the system logic when it is created:

\begin{cppcode}
    constexpr std::uint8_t ledPin{2U};
    constexpr std::uint8_t buttonPin{3U};

    // Create ESP32-S3 factory.
    driver::factory::Esp32s3 factory{};

    // Create and run system logic.
    app::logic::Logic logic{factory, ledPin, buttonPin};
    logic.run();
\end{cppcode}

\noindent This moves the creation of drivers away from the system logic. Instead, the factory is
responsible for selecting and creating the correct concrete implementation. The system logic can
therefore operate on interfaces without knowing the underlying hardware.

\subsection{Dynamic memory allocation in the factory}
\label{c4:sec:dynamic}
When instances are created in this way, they normally need to be allocated dynamically, unless the
factory instead returns a reference to a singleton class.

This can be done using either:
\begin{itemize}
  \item Raw pointers (\code{new} and \code{delete}).
  \item Smart pointers (\code{std::make_unique()}).
\end{itemize}
Instead of allocating an instance of \code{driver::gpio::Esp32s3} dynamically ourselves:

\begin{cppcode}
auto* gpio = new driver::gpio::Esp32s3(pin);
\end{cppcode}

\noindent we let a factory perform this task:

\begin{cppcode}
auto* gpio = factory.gpio(pin);
\end{cppcode}

\noindent This means that:
\begin{itemize}
  \item The system logic does not know which concrete class is used.
  \item It only knows that it receives an object that satisfies a certain interface.
\end{itemize}

\subsection{Why do we use a factory?}
\label{c4:sec:whyfactory}
In many embedded systems, a factory plays a role similar to a hardware abstraction layer (HAL) or a
board support package (BSP); it centralizes the creation of hardware-specific objects.

In an embedded system we often want to be able to:

\begin{narrowtable}{@{}ll@{}}
  \thead{Situation} & \thead{What we want to do} \\
  \midrule
  Real hardware & Use real drivers \\
  Testing/simulation & Use stubs \\
  New platform & Replace drivers \\
\end{narrowtable}

\noindent If the system logic creates drivers itself:

\begin{cppcode}
auto* led = new driver::gpio::Esp32s3(pin);
\end{cppcode}

\noindent then:
\begin{itemize}
  \item The logic is tied to ESP32-S3.
  \item The system cannot be tested without hardware.
\end{itemize}
With a factory:

\begin{cppcode}
auto* led = factory.gpio(pin);
\end{cppcode}

\noindent we can switch between:

\begin{cppcode}
driver::factory::Esp32s3 factory{};
\end{cppcode}

\noindent and:

\begin{cppcode}
driver::factory::Stub factory{};
\end{cppcode}

\noindent without changing the system logic.

This is called \emph{dependency injection via factory}.

\subsection{Overview of the architecture in a typical embedded system}
\label{c4:sec:architecture}
Assume that we have an embedded system where the software is divided into four layers:

\begin{console}
app::logic::Logic
        |
        v
driver::factory::Interface
        |
        +------------+
        |            |
        v            v
factory::Esp32s3   factory::Stub
        |            |
        v            v
gpio::Esp32s3      gpio::Stub
\end{console}

\noindent The system logic therefore does not know whether it is running on real hardware or on
stubs. It only knows about the interfaces.

\subsection{Example implementation using a factory}
\label{c4:sec:rawexample}

\subsubsection{Step 1: Interface for GPIO}
We create a simple GPIO interface in a file \file{driver/gpio/interface.hpp}, as shown below:

\begin{cppcode}
#pragma once

namespace driver::gpio
{
class Interface
{
public:
    virtual ~Interface() noexcept = default;
    [[nodiscard]] virtual bool read() const noexcept = 0;
    virtual void write(bool state) noexcept = 0;
    virtual void toggle() noexcept = 0;
};
} // namespace driver::gpio
\end{cppcode}

\subsubsection{Step 2: Real driver}
We create a GPIO driver for the ESP32-S3 in a file \file{driver/gpio/esp32s3.hpp}, as shown below:

\begin{cppcode}
#pragma once

#include <cstdint>

#include "driver/gpio/interface.hpp"

namespace driver::gpio
{
class Esp32s3 final : public Interface
{
public:
    explicit Esp32s3(const std::uint8_t pin) noexcept
        : myPin{pin}
    {
        // ESP-specific initialization here!
    }

    ~Esp32s3() noexcept override
    {
        // ESP-specific cleanup here!
    }

    [[nodiscard]] bool read() const noexcept override
    {
        // ESP-specific code here!
        return false;
    }

    void write(const bool state) noexcept override
    {
        // ESP-specific code here!
        (void) (state);
    }

    void toggle() noexcept override
    {
        // ESP-specific code here!
    }

    Esp32s3()                          = delete;
    Esp32s3(const Esp32s3&)            = delete;
    Esp32s3(Esp32s3&&)                 = delete;
    Esp32s3& operator=(const Esp32s3&) = delete;
    Esp32s3& operator=(Esp32s3&&)      = delete;

private:
    const std::uint8_t myPin;
};
} // namespace driver::gpio
\end{cppcode}

\subsubsection{Step 3: Stub}
We create a GPIO stub in a file \file{driver/gpio/stub.hpp}, as shown below:

\begin{cppcode}
#pragma once

#include <cstdint>

#include "driver/gpio/interface.hpp"

namespace driver::gpio
{
class Stub final : public Interface
{
public:
    Stub() noexcept
        : myState{false}
    {}
    ~Stub() noexcept override = default;

    [[nodiscard]] bool read() const noexcept override { return myState; }
    void write(const bool state) noexcept override { myState = state; }
    void toggle() noexcept override { myState = !myState; }

    Stub(const Stub&)            = delete;
    Stub(Stub&&)                 = delete;
    Stub& operator=(const Stub&) = delete;
    Stub& operator=(Stub&&)      = delete;

private:
    bool myState;
};
} // namespace driver::gpio
\end{cppcode}

\subsubsection{Step 4: Factory interface}
We create a factory interface in a file \file{driver/factory/interface.hpp} in order to implement
separate factories:
\begin{itemize}
  \item A factory for creating instances for ESP32-S3.
  \item A factory for creating stub instances.
\end{itemize}

\begin{cppcode}
#pragma once

#include <cstdint>

namespace driver
{
/** GPIO driver interface. */
namespace gpio { class Interface; }
} // namespace driver

namespace driver::factory
{
class Interface
{
public:
    virtual ~Interface() noexcept = default;
    [[nodiscard]] virtual gpio::Interface* gpio(std::uint8_t pin) noexcept = 0;
};
} // namespace driver::factory
\end{cppcode}

\noindent Note that \code{gpio()} is marked \code{[[nodiscard]]}: it allocates a new object with
\code{new} and hands ownership to the caller. Discarding the returned pointer would lose the only
reference to that object, leaking it --- there would be no way to \code{delete} it later.

\subsubsection{Step 5: Real factory}
We then create an ESP32-S3 factory in a file \file{driver/factory/esp32s3.hpp} in order to
construct ESP32 instances:

\begin{cppcode}
#pragma once

#include <cstdint>

#include "driver/factory/interface.hpp"
#include "driver/gpio/esp32s3.hpp"

namespace driver::factory
{
class Esp32s3 final : public Interface
{
public:
    Esp32s3() noexcept = default;
    ~Esp32s3() noexcept override = default;

    [[nodiscard]] gpio::Interface* gpio(const std::uint8_t pin) noexcept override
    {
        return new gpio::Esp32s3(pin);
    }

    Esp32s3(const Esp32s3&)            = delete;
    Esp32s3(Esp32s3&&)                 = delete;
    Esp32s3& operator=(const Esp32s3&) = delete;
    Esp32s3& operator=(Esp32s3&&)      = delete;
};
} // namespace driver::factory
\end{cppcode}

\subsubsection{Step 6: Stub factory}
We also create a stub factory in a file \file{driver/factory/stub.hpp} in order to construct stub
instances:

\begin{cppcode}
#pragma once

#include <cstdint>

#include "driver/factory/interface.hpp"
#include "driver/gpio/stub.hpp"

namespace driver::factory
{
class Stub final : public Interface
{
public:
    Stub() noexcept = default;
    ~Stub() noexcept override = default;

    [[nodiscard]] gpio::Interface* gpio(const std::uint8_t pin) noexcept override
    {
        // Ignore the pin number since it is not used by the stub.
        (void) (pin);
        return new gpio::Stub();
    }

    Stub(const Stub&)            = delete;
    Stub(Stub&&)                 = delete;
    Stub& operator=(const Stub&) = delete;
    Stub& operator=(Stub&&)      = delete;
};
} // namespace driver::factory
\end{cppcode}

\subsubsection{Step 7: System logic using a factory}
We then create a logic class that uses a given factory to create instances, in a file
\file{app/logic/logic.hpp}:
\begin{itemize}
  \item In the constructor we allocate memory for the GPIO instances \code{myLed} and
    \code{myButton} via the factory.
  \item In the destructor we free the allocated resources, i.e.\ the GPIO instances --- since we
    have allocated memory using raw pointers we are responsible for doing this ourselves.
  \item In the method \code{run()}, \code{myLed} is toggled on the rising edge of the push button.
  \item Because \code{myLed} and \code{myButton} are pointers, the arrow operator \code{->} is used
    to call their methods.
  \item The copy and move operations are deleted. This is particularly important here, because
    \code{Logic} owns two raw pointers and deletes them in its destructor. Had we not deleted them,
    the compiler would have generated a copy constructor that copies the two \emph{pointer values},
    so a copy would refer to the exact same two GPIO objects. Both objects would then delete them
    when destroyed, which means each GPIO object would be deleted twice; a \textbf{double
    deletion}, which is undefined behavior, would occur. The same applies to the copy assignment
    operator.
\end{itemize}

\note This is a general rule, and it is sometimes referred to as the \emph{rule of three}: a class
that needs a destructor almost always needs its copy operations written or deleted as well. For a
class such as \code{Logic}, which represents one system with one set of drivers, deleting them is
the right choice.

In this example, the system logic owns the driver objects. The pointers returned by the factory are
therefore owned by \code{Logic}, which means that the system logic is also responsible for deleting
them in the destructor:

\begin{cppcode}
#pragma once

#include <cstdint>

#include "driver/factory/interface.hpp"
#include "driver/gpio/interface.hpp"

namespace app::logic
{
class Logic final
{
public:
    explicit Logic(driver::factory::Interface& factory,
                   const std::uint8_t ledPin, const std::uint8_t buttonPin) noexcept
        : myLed{factory.gpio(ledPin)}
        , myButton{factory.gpio(buttonPin)}
    {}

    ~Logic() noexcept
    {
        delete myLed;
        delete myButton;
    }

    void run() noexcept
    {
        bool buttonPrev{false};

        while (1)
        {
            // Read the current button input.
            const bool buttonCurrent{myButton->read()};

            // Toggle the LED on pressdown (rising edge).
            if (buttonCurrent && !buttonPrev)
            {
                myLed->toggle();
            }
            // Store the current button input for next comparison.
            buttonPrev = buttonCurrent;
        }
    }

    // Deleting the copy operations is essential here: a copy would hold the same two pointers,
    // and both objects would delete the same GPIO instances (double deletion).
    Logic()                        = delete; // No default constructor.
    Logic(const Logic&)            = delete; // No copy constructor.
    Logic(Logic&&)                 = delete; // No move constructor.
    Logic& operator=(const Logic&) = delete; // No copy assignment.
    Logic& operator=(Logic&&)      = delete; // No move assignment.

private:
    driver::gpio::Interface* myLed;
    driver::gpio::Interface* myButton;
};
} // namespace app::logic
\end{cppcode}

\subsubsection{Step 8: \code{main}}
In the function \code{main}, we create a factory and pass it to the system logic to create the
objects.

Below is an example of how a factory could be used to run the system logic on an ESP32-S3 by
passing an instance of \code{driver::factory::Esp32s3} to the system logic, followed by calling the
method \code{run()}:

\begin{cppcode}
#include <cstdint>

#include "driver/factory/esp32s3.hpp"
#include "app/logic/logic.hpp"

int main()
{
    constexpr std::uint8_t ledPin{2U};
    constexpr std::uint8_t buttonPin{3U};

    // Create system logic and initialize the system.
    driver::factory::Esp32s3 factory{};
    app::logic::Logic logic{factory, ledPin, buttonPin};

    // Run the system continuously.
    logic.run();
    return 0;
}
\end{cppcode}

\noindent To run with stub drivers, it is sufficient to switch to the stub factory
\code{driver::factory::Stub}, as shown below. This is well suited for testing and will be
demonstrated in a later QAcademy course on software testing:

\begin{cppcode}
#include <cstdint>

#include "driver/factory/stub.hpp"
#include "app/logic/logic.hpp"

int main()
{
    constexpr std::uint8_t ledPin{2U};
    constexpr std::uint8_t buttonPin{3U};

    // Create system logic and initialize the system.
    driver::factory::Stub factory{};
    app::logic::Logic logic{factory, ledPin, buttonPin};

    // Run the system continuously.
    logic.run();
    return 0;
}
\end{cppcode}
